% META: {"id": "1SPE-VARALEA-CR-X3", "chapitre": "1SPE-VARIABLES-ALEATOIRES", "type_objet": "cours", "capacites_codes": [], "extension_codes": ["X3"], "programme_alignment": "OPTIONAL_EXTENSION", "extension_label": "Approfondissement — Vers la Terminale", "sources_inspiration": [], "status": "generated"}

\begin{center}\textbf{Approfondissement — Vers la Terminale}\end{center}

\section{Variance et écart type après une transformation affine}

\propriete{
Si $Y=aX+b$, où $a$ et $b$ sont deux réels, alors
\[
V(Y)=a^2V(X),\qquad \sigma(Y)=|a|\sigma(X).
\]
Ajouter une constante déplace toutes les valeurs sans modifier leur
dispersion ; multiplier par $a$ multiplie les écarts à la moyenne par $|a|$.
}

\exemple{
Une recette $X$, exprimée en euros, vérifie $V(X)=5$. Après des frais fixes
de $3$ euros, le bénéfice est $Y=X-3$, donc $V(Y)=5$. Si toutes les sommes
sont ensuite doublées, pour $Z=2Y$ on obtient $V(Z)=4V(Y)=20$.
}

\erreurFrequente{
La variance n'est pas affine : $V(aX+b)$ ne vaut ni $aV(X)+b$ ni $aV(X)$.
Le coefficient est mis au carré et la constante $b$ disparaît.
}
